% from ega2-1-fr.tex
\oldpage[II]{111}

Nous démontrerons d'abord le lemme suivant~:

\begin{lemma}[6.1.4.1]
\label{II.6.1.4.1-fr}
Soient $A$ un anneau, $M$ un $A$-module, $(g_i)_{1\leq i\leq m}$ un système
fini d'éléments de $A$ tels que les $D(g_i)$ ($1\leq i\leq m$) recouvrent
$\operatorname{Spec}(A)$. Si, pour tout $i$, $M_{g_i}$ est un
$A_{g_i}$-module de type fini, alors $M$ est un $A$-module de type fini.
\end{lemma}

On peut en effet supposer que $M_{g_i}$ admet un système de générateurs fini
$(m_{ij}/g_i^n)$ avec $m_{ij}\in M$, $n$ étant le même pour tous les indices
$i$. Montrons que les $m_{ij}$ forment un système de générateurs de $M$. Soit
$M'$ le sous-$A$-module de $M$ engendré par ces éléments, et soit $m$ un
élément de $M$. Par hypothèse, pour chaque $i$, il existe des $a_{ij}\in A$ et
un entier $p$ (indépendant de $i$) tels que, dans $M_{g_i}$,
$m/1=(\sum_j a_{ij}m_{ij})/g_i^p$~; cela entraîne qu'il existe un entier
$r\geq p$ tel que, pour tout $i$, on ait $g_i^r m\in M'$. Or, comme les
$D(g_i^r)=D(g_i)$ recouvrent $\operatorname{Spec}(A)$, l'idéal de $A$
engendré par les $g_i^r$ est égal à $A$, autrement dit il existe des éléments
$a_i\in A$ tels que $\sum_i a_i g_i^r=1$~; par suite
$m=(\sum_i a_i g_i^r)m\in M'$, d'où le lemme.

Cela étant, on sait déjà (\hyperref[II.1.3.2-fr]{1.3.2}) que $f^{-1}(U)$ est
affine. Si $\varphi$ est l'homomorphisme $A\to B$ correspondant à $f$, il
existe un recouvrement fini de $U$ par des ouverts $D(g_i)$ ($g_i\in A$) tels
que, si $h_i=\varphi(g_i)$, $B_{h_i}$ soit une algèbre entière (resp. entière
finie) sur $A_{g_i}$. En effet, il y a un recouvrement de $U$ par des ouverts
affines $V_\alpha\subset U$ tels que si
$A_\alpha=A(V_\alpha)$, $B_\alpha=A(f^{-1}(V_\alpha))$ soit une algèbre
entière (resp. finie) sur $A_\alpha$. Tout $x\in U$ appartient à un
$V_\alpha$, donc il existe $g\in A$ tel que $x\in D(g)\subset V_\alpha$~; si
$g_\alpha$ est l'image de $g$ dans $A_\alpha$, on a
$A(D(g))=A_g=(A_\alpha)_{g_\alpha}$~; soit $h=\varphi(g)$, et soit $h_\alpha$
l'image de $g_\alpha$ dans $B_\alpha$~; on a
\[
A(D(h))=B_h=(B_\alpha)_{h_\alpha}
\]
et comme $B_\alpha$ est entière (resp. finie) sur $A_\alpha$,
$(B_\alpha)_{h_\alpha}$ est entière (resp. finie) sur
$(A_\alpha)_{g_\alpha}$. Il suffit maintenant d'utiliser le fait que $U$ est
quasi-compact pour obtenir le recouvrement cherché.

Si on suppose d'abord que les $B_{h_i}$ sont entières finies sur les
$A_{g_i}$, comme, en tant que $A_{g_i}$-module, $B_{h_i}$ s'écrit aussi
$B_{g_i}$, le lemme (\hyperref[II.6.1.4.1-fr]{6.1.4.1}) montre que dans ce cas
$B$ est un $A$-module de type fini.

Supposons seulement maintenant chaque $B_{h_i}$ entière sur $A_{g_i}$~; soit
$b\in B$, et soit $C$ la sous-$A$-algèbre de $B$ engendrée par $b$. Pour tout
$i$, $C_{h_i}$ est l'algèbre sur $A_{g_i}$ engendrée par $b/1$ dans
$B_{h_i}$~; il résulte de l'hypothèse que chaque $C_{h_i}$ est un
$A_{g_i}$-module de type fini, donc
(\hyperref[II.6.1.4.1-fr]{6.1.4.1}) $C$ est un $A$-module de type fini, ce qui
prouve que $B$ est entière sur $A$.

\begin{proposition}[6.1.5]
\label{II.6.1.5-fr}
\begin{enumerate}
\item[(i)] Une immersion fermée est finie (et a fortiori entière).
\item[(ii)] Le composé de deux morphismes finis (resp. entiers) est fini
(resp. entier).
\item[(iii)] Si $f:X\to Y$ est un $S$-morphisme fini (resp. entier),
$f_{(S')}:X_{(S')}\to Y_{(S')}$ est fini (resp. entier) pour toute extension
$S'\to S$ de la base.
\item[(iv)] Si $f:X\to Y$ et $g:X'\to Y'$ sont deux $S$-morphismes finis
(resp. entiers), il en est de même de
$f\times_S g:X\times_S X'\to Y\times_S Y'$.
\item[(v)] Si $f:X\to Y$ et $g:Y\to Z$ sont deux morphismes tels que $g\circ
f$ soit fini (resp. entier) et $g$ séparé, alors $f$ est fini (resp. entier).
\item[(vi)] Si $f:X\to Y$ est un morphisme fini (resp. entier), il en est de
même de $f_{\mathrm{red}}$.
\end{enumerate}
\end{proposition}

